\chapter{مدیریت عملیات چاپ}

پس از گذراندن چندین فصل پیرامون پردازش و دستکاری متن، اکنون زمان آن فرا رسیده است که نتایج حاصل را بر روی کاغذ منتقل کنیم. در این فصل، ابزارهای مبتنی بر رابط خط فرمان را بررسی خواهیم کرد که جهت چاپ فایل‌ها و مدیریت عملکرد چاپگر طراحی شده‌اند. از آنجایی که پیکربندی سیستم چاپ در توزیع‌های مختلف لینوکس متفاوت است و معمولاً در حین نصب سیستم‌عامل به‌طور خودکار انجام می‌شود، در این بخش به تنظیمات سخت‌افزاری نخواهیم پرداخت. توجه داشته باشید که جهت اجرای تمرینات این فصل، برخورداری از یک سرویس چاپ فعال و پیکربندی‌شده الزامی است.

در این فصل، دستورات زیر را تحلیل خواهیم کرد:

\begin{itemize}
  \item \cmd{pr}: آماده‌سازی و تبدیل فایل‌های متنی برای چاپ.
  \item \cmd{lp/lpr}: ارسال فایل‌ها به چاپگر.
  \item \cmd{a2ps}: قالب‌بندی فایل‌ها جهت خروجی در چاپگرهای \en{PostScript}.
  \item \cmd{lpstat}: نمایش وضعیت جاری چاپگرها.
  \item \cmd{lpq}: بررسی وضعیت صف انتظار چاپ (\en{Print Queue}).
  \item \cmd{lprm}: لغو و حذف وظایف (\en{Jobs}) از صف چاپ.
\end{itemize}

\section{تاریخچه چاپ}

برای درک عمیق قابلیت‌های سیستم چاپ در سیستم‌عامل‌های شبه‌یونیکس، بازخوانی تاریخچه‌ی آن ضروری است. ریشه‌های چاپ در این محیط‌ها به نخستین روزهای شکل‌گیری سیستم‌عامل بازمی‌گردد؛ زمانی که ماهیت چاپگرها و الگوهای کاربری آن‌ها تفاوت بنیادینی با استانداردهای امروزی داشت.

\subsection{چاپ در عصر رایانه‌های متمرکز}

در دوران پیش از ظهور رایانه‌های شخصی (\en{PC})، چاپگرها همانند خودِ رایانه‌ها، دستگاه‌هایی عظیم، گران‌قیمت و متمرکز بودند. یک کاربر معمولی در سال ۱۹۸۰، فرآیندهای خود را از طریق ترمینالی پیش می‌برد که به یک رایانه‌ی مرکزی در فاصله‌ای دور متصل بود. چاپگرها نیز معمولاً در مجاورت واحد مرکزی مستقر بوده و توسط اپراتورهای متخصص مدیریت می‌شدند.

به دلیل هزینه‌ی گزاف و تمرکز سخت‌افزاری، چندین کاربر به‌طور مشترک از یک چاپگر بهره می‌بردند. برای تفکیک خروجی‌های هر فرد، سازوکاری ابداع شد که در آن یک «صفحه‌ی شناسایی» (\en{Banner Page}) حاوی نام کاربر، در ابتدای هر دستور چاپ درج می‌شد؛ در نهایت، تیم پشتیبانی خروجی‌های چاپی هر روز را دسته‌بندی کرده و به کاربران مربوطه تحویل می‌داد.

\subsection{چاپگرهای مبتنی بر نویسه}

فناوری چاپ در دهه‌ی ۸۰ میلادی از دو منظر کلیدی با تکنولوژی معاصر متفاوت بود. نخست آنکه چاپگرهای آن دوره غالباً از نوع «ضربه‌ای» (\en{Impact Printer}) بودند. این دستگاه‌ها با استفاده از مکانیزم‌های مکانیکی، نوار آغشته به جوهر را بر روی کاغذ می‌کوبیدند تا اثر نویسه شکل بگیرد. دو فناوری پیشرو در آن زمان، چاپگرهای چرخ آفتابگردان (\en{Daisy-Wheel}) و «ماتریس‌نقطه‌ای» (\en{Dot-Matrix}) بودند.

دوم و مهم‌تر از آن، این چاپگرها از مجموعه‌ای ثابت از نویسه‌ها (\en{Character Sets}) بهره می‌بردند که به‌صورت سخت‌افزاری در دل دستگاه تعبیه شده بود. برای مثال، یک چاپگر چرخ آفتابگردان (\en{Daisy-Wheel}) صرفاً قادر به درج علائمی بود که به‌صورت فیزیکی بر روی پره‌های چرخ آن حک شده بودند. این ویژگی، چاپگرها را به نسخه‌های صنعتی و پرسرعت ماشین‌تحریر شبیه می‌کرد. مشابه اکثر ماشین‌های تحریر، فرآیند چاپ با استفاده از قلم‌های تک‌فاصله (\en{Monospaced} یا \en{Fixed-Width}) انجام می‌شد؛ بدین معنا که هر نویسه، صرف‌نظر از شکل ظاهری، فضای افقی یکسانی را اشغال می‌کرد.

اکثر چاپگرها در آن زمان با تراکم ۱۰ نویسه در اینچ (\en{CPI}) به‌صورت افقی و ۶ سطر در اینچ (\en{LPI}) به‌صورت عمودی فعالیت می‌کردند. بر اساس این استاندارد، یک کاغذ در ابعاد \en{US Letter} دارای ظرفیت ۸۵ نویسه در عرض و ۶۶ سطر در ارتفاع بود. با کسر حاشیه امن (\en{Safety Margins})، عدد ۸۰ نویسه به‌عنوان حداکثر عرض استاندارد برای هر سطر تثبیت شد. این موضوع به‌خوبی تبیین می‌کند که چرا عرض نمایشگرهای ترمینال (و شبیه‌سازهای مدرن آن) به‌طور پیش‌فرض ۸۰ نویسه است. در واقع، استفاده از قلم تک‌فاصله در ترمینالی با عرض ۸۰ نویسه، نخستین تجربه‌ی بشری از مفهوم «آنچه می‌بینی، همان است که چاپ می‌شود» (\en{WYSIWYG - What You See Is What You Get}) را رقم زد.

انتقال داده‌ها به این چاپگرها در قالب جریانی ساده از بایت‌ها صورت می‌گرفت؛ بایت‌هایی که مستقیماً کدهای اسکی (\en{ASCII}) نویسه‌های مورد نظر را حمل می‌کردند. علاوه بر نویسه‌های متنی، کدهای کنترلی سطح پایین \en{ASCII} نیز برای هدایت مکانیکی دستگاه ارسال می‌شدند؛ دستوراتی نظیر «بازگشت به ابتدای سطر» (\en{Carriage Return})، «خط جدید» (\en{Line Feed}) و «انتقال کاغذ به ابتدای صفحه بعدی» (\en{Form Feed}). با ترکیب هوشمندانه‌ی این کدهای کنترلی، حتی ایجاد جلوه‌های بصری محدودی مانند «برجسته کردن متن» (\en{Boldface}) نیز میسر بود؛ این کار از طریق چاپ یک نویسه، بازگشت به عقب (\en{Backspace}) و چاپ مجدد همان نویسه روی اثر قبلی انجام می‌شد تا ضخامت جوهر روی کاغذ افزایش یابد. در واقع می‌توان این فرآیند را از نزدیک مشاهده کرد؛ کافی است با استفاده از ابزار \cmd{nroff} یک صفحه راهنما (\en{man page}) را رندر نموده و سپس خروجی حاصل را جهت بررسی نویسه‌های کنترلی به دستور \cmd{cat -A} هدایت کنیم:

\begin{latin}
\begin{lstlisting}
[me@linuxbox ~]$ zcat /usr/share/man/man1/ls.1.gz | nroff -man | cat -A | head
LS(1)                     User Commands                           LS(1)
$
$
$
N^HNA^HAM^HME^HE$
     ls - list directory contents$
$
S^HSY^HYN^HNO^HOP^HPS^HSI^HIS^HS$
     l^Hls^Hs [_^HO_^HP_^HT_^HI_^HO_^HN]... [_^HF_^HI_^HL_^HE]...$
\end{lstlisting}
\end{latin}

در خروجی فوق، نویسه‌ی \cmd{\^{}H} (معادل ترکیبی \cmd{Ctrl+h}) بیانگر کاراکتر پس‌بَر یا همان \en{Backspace} است. این تکنیک برای ایجاد جلوه‌ی پررنگ (\en{Bold}) از طریق چاپ مجدد یک نویسه روی خودش به کار می‌رفته است. همچنین در بخش‌های دیگر، توالی نویسه‌های \en{Backspace} و زیرخط (\en{Underscore}) برای تولید متن زیرخط‌دار (\en{Underlined}) قابل مشاهده است.

\subsection{چاپگرهای گرافیکی}

تحول در فناوری چاپگرهای گرافیکی و توسعهٔ رابط‌های کاربری گرافیکی (\en{GUI}) دگرگونی‌های بنیادینی را در فناوری چاپ رقم زد. همگام با گذار رایانه‌ها به سمت محیط‌های بصری و مبتنی بر تصویر، فرآیند چاپ نیز از روش‌های نویسه‌محور (\en{Character-based}) به تکنیک‌های گرافیکی تغییر یافت. این گذار با ظهور چاپگرهای لیزری مقرون‌به‌صرفه شتاب گرفت؛ دستگاه‌هایی که به‌جای چاپ نویسه‌های ثابت فیزیکی، قادر بودند نقاط بسیار ریزی را در هر مختصاتی از سطح کاغذ ترسیم کنند. این نوآوری، امکان استفاده از قلم‌های متناسب (\en{Proportional Fonts}) – مشابه آنچه در صنعت نشر به کار می‌رود – و همچنین چاپ تصاویر و نمودارهای با وضوح بالا را میسر ساخت.

با این حال، انتقال از ساختار متنی به گرافیکی با یک چالش فنی جدی همراه بود. دلیل این امر در تفاوت فاحش حجم داده‌های انتقالی نهفته است: حجم دادهٔ مورد نیاز برای پر کردن یک صفحه در چاپگرهای نویسه‌محور (با فرض ۶۰ سطر و ۸۰ نویسه در هر سطر) تنها ۴۸۰۰ بایت است:

\begin{latin}
\begin{lstlisting}
(60 x 80) = 4,800 bytes
\end{lstlisting}
\end{latin}

در مقابل، یک چاپگر لیزری با وضوح ۳۰۰ نقطه در اینچ (\en{DPI}) در یک ناحیهٔ چاپی ۸ در ۱۰ اینچ، به حجمی معادل ۹۰۰,۰۰۰ بایت نیاز دارد:

\begin{latin}
\begin{lstlisting}
(8 x 300) x (10 x 300) / 8 = 900,000 bytes
\end{lstlisting}
\end{latin}

بسیاری از شبکه‌های رایانه‌ای اولیه قادر به انتقال سریع حدود یک مگابایت داده برای هر صفحه نبودند؛ بنابراین، ابداع راهکاری هوشمندانه ضروری می‌نمود.

این راهکار، زبان توصیف صفحه (\en{Page Description Language – PDL}) بود. \en{PDL} در واقع یک زبان برنامه‌نویسی تخصصی برای تبیین محتوای بصری صفحه است. این زبان به‌جای ارسال نقشهٔ نقاط (\en{Bitmap})، دستورالعمل‌هایی صادر می‌کند: «به مختصات مشخص برو، نویسهٔ '\en{a}' را با قلم \en{Helvetica} و اندازهٔ ۱۰ درج کن و سپس به موقعیت بعدی برو». برجسته‌ترین نمونه در این حوزه، زبان \en{PostScript} از شرکت \en{Adobe Systems} است که همچنان جایگاه خود را در صنعت حفظ کرده است. \en{PostScript} یک زبان برنامه‌نویسی کامل برای تایپوگرافی و تصویرسازی گرافیکی است که به‌صورت پیش‌فرض از ۳۵ قلم استاندارد با کیفیت بالا پشتیبانی می‌کند و قابلیت بارگذاری قلم‌های جدید را نیز داراست.

در ابتدا، مفسر \en{PostScript} مستقیماً درون سخت‌افزار چاپگر تعبیه می‌شد. این رویکرد معضل پهنای باند شبکه را مرتفع کرد؛ چرا که فایل دستورالعمل‌های \en{PostScript}، اگرچه حجیم‌تر از متن ساده بود، اما بسیار سبک‌تر از نقشهٔ بیتی کامل صفحه (\en{Bitmap}) عمل می‌کرد. یک چاپگر مجهز به \en{PostScript}، ورودی را دریافت کرده و با بهره‌گیری از پردازنده و حافظهٔ داخلی خود (که گاه قدرتمندتر از رایانهٔ میزبان بود)، مفسر مخصوصی را برای رندر کردن صفحه و تولید الگوی نقاط نهایی اجرا می‌کرد. این فرآیند تبدیل دستورات به نقشهٔ نقاط، «پردازندهٔ تصویر شطرنجی» (\en{Raster Image Processor – RIP}) نامیده می‌شود.

با افزایش توان پردازشی رایانه‌ها و سرعت شبکه‌ها، فرآیند \en{RIP} از حافظهٔ داخلی چاپگر به سیستم میزبان انتقال یافت. این جابه‌جایی باعث شد چاپگرهای با کیفیت بالا با قیمت بسیار کمتری به بازار عرضه شوند. امروزه، در حالی که بسیاری از چاپگرهای ارزان‌قیمت (\en{Host-based}) کاملاً به پردازش رایانهٔ میزبان وابسته‌اند، چاپگرهای پیشرفتهٔ \en{PostScript} همچنان در محیط‌های حرفه‌ای مورد استفاده قرار می‌گیرند.

\section{سیستم چاپ در لینوکس}

توزیع‌های مدرن لینوکس برای مدیریت و اجرای فرآیندهای چاپ، از دو مجموعه‌ی نرم‌افزاری کلیدی بهره می‌برند. نخستین مجموعه، \en{CUPS} (مخفف \en{Common Unix Printing System}) یا «سامانه چاپ مشترک یونیکس» است که وظیفه‌ی تأمین راه‌اندازهای چاپگر و مدیریت صف‌های عملیاتی را بر عهده دارد. مجموعه‌ی دوم، \en{Ghostscript} نام دارد که به عنوان یک مفسر نرم‌افزاری برای زبان \en{PostScript} عمل کرده و نقش پردازندهٔ تصویر شطرنجی (\en{RIP}) را ایفا می‌کند.

سیستم \en{CUPS} مسئولیت تعریف و نگهداری صف‌های چاپ (\en{Print Queues}) را بر عهده دارد. همان‌گونه که در پیشینه‌ی تاریخی این فناوری اشاره شد، معماری چاپ در محیط‌های یونیکسی اساساً برای مدیریت یک سخت‌افزار متمرکز در میان چندین کاربر طراحی شده است. از آنجا که سرعت پردازش و خروجی چاپگرها به‌طور ذاتی بسیار کمتر از توان محاسباتی رایانه‌هاست، وجود سازوکاری جهت زمان‌بندی (\en{Scheduling}) وظایف و سازمان‌دهی اولویت‌ها ضروری است. علاوه بر این، \en{CUPS} قادر است انواع فایل‌های ورودی (از جمله متن ساده، \en{PDF} و تصاویر) را به‌طور خودکار شناسایی کرده و پیش از ارسال به چاپگر، آنها را به قالبی مناسب و قابل چاپ برای دستگاه مقصد تبدیل نماید.

\section{آماده‌سازی اسناد برای چاپ}

به‌عنوان کاربران رابط خط فرمان، تمرکز اصلی ما غالباً بر چاپ مستندات متنی معطوف است؛ هرچند که چاپ سایر قالب‌های فایلی نیز به‌طور کامل در لینوکس پشتیبانی می‌شود.

\subsection{دستور pr: آماده‌سازی متن برای چاپ}

ما در فصل گذشته نگاهی گذرا به ابزار \cmd{pr} داشتیم. در این بخش، برخی از گزینه‌های متعدد این دستور را که به‌طور اختصاصی در فرآیند چاپ کاربرد دارند، تحلیل خواهیم کرد. همان‌طور که در مرور تاریخی اشاره شد، چاپگرهای نویسه‌محور قدیمی از قلم‌های تک‌فاصله (\en{Monospaced}) بهره می‌بردند که موجب می‌شد هر سطر و هر صفحه دارای تعداد ثابتی از کاراکترها باشد. ابزار \cmd{pr} وظیفه دارد متن را به‌گونه‌ای تنظیم کند که در ابعاد یک صفحه‌ی استاندارد، همراه با سرآیند‌ها (\en{Headers}) و حاشیه‌های (\en{Margins}) انتخابی، به‌درستی جای بگیرد.

در جدول زیر، کاربردی‌ترین گزینه‌های دستور \cmd{pr} برای مدیریت چیدمان متن خلاصه شده است:

\begin{table}[htbp]
\centering
\caption{گزینه‌های متداول دستور \cmd{pr}}
\label{tab:pr-options}
\begin{tabularx}{\textwidth}{l X}
\toprule
\textbf{گزینه} & \textbf{شرح عملکرد} \\
\midrule
\cmd{+first[:last]} & چاپ بازه‌ای از صفحات؛ از شماره‌ی صفحه‌ی \en{first} تا شماره‌ی صفحه‌ی اختیاریِ \en{last}. \\
\addlinespace
\cmd{-columns} & سازمان‌دهیِ محتوای متنی در تعداد ستونِ مشخص‌شده (\en{columns}). \\
\addlinespace
\cmd{-a} & تغییر جهت چیدمانِ خروجیِ چندستونه از حالت عمودی (پیش‌فرض) به حالت افقی؛ بدین ترتیب، محتوا ابتدا در طول هر سطر و سپس در ستون بعدی پر می‌شود. \\
\addlinespace
\cmd{-d} & اعمال فاصله‌گذاری دوبرابر (\en{Double-space}) میان سطرهای خروجی. \\
\addlinespace
\cmd{-D ``format''} & شخصی‌سازیِ قالب نمایش تاریخ در سرآیند (\en{Header})، با استفاده از همان نشانه‌های قالب‌بندیِ رایج در دستور \cmd{date}. \\
\addlinespace
\cmd{-f} & استفاده از نویسه‌ی کنترلیِ \en{Form Feed} برای جداسازیِ صفحات، به‌جای درج ردیفی از سطرهای خالی. \\
\addlinespace
\cmd{-h ``header''} & جایگزینیِ نام فایل با متنی دلخواه، در بخش میانیِ سرآیند هر صفحه. \\
\addlinespace
\cmd{-l length} & تعیین طول صفحه بر حسب تعداد سطر (\en{length}). مقدار پیش‌فرض این گزینه ۶۶ سطر است. \\
\addlinespace
\cmd{-n} & افزودن شماره‌گذاریِ خودکار به ابتدای هر سطر از خروجی. \\
\addlinespace
\cmd{-o offset} & ایجاد حاشیه‌ای در سمت چپ صفحه، به اندازه‌ی \en{offset} نویسه‌ی فاصله. \\
\addlinespace
\cmd{-w width} & تعیین عرض کل صفحه بر حسب تعداد نویسه (\en{width}). مقدار پیش‌فرض این گزینه ۷۲ نویسه است. \\
\bottomrule
\end{tabularx}
\end{table}

دستور \cmd{pr} اغلب در نقش یک فیلتر در پایپ‌لاین (\en{Pipeline}) برای آرایش محتوا به کار می‌رود. در مثال زیر، فهرست محتویات دایرکتوری \file{/usr/bin/} را استخراج کرده و آن را به خروجی صفحه‌بندی‌شده در سه ستون تبدیل می‌کنیم:

\begin{latin}
\begin{lstlisting}
[me@linuxbox ~]$ ls /usr/bin | pr -3 -w 65 | head

2025-02-18 14:00                                      Page 1

[                 apturl                    bsd-write
411toppm          ar                        bsh
a2p               arecord                   btcf lash
a2ps              arecordmidi               bug-buddy
a2ps-lpr-wrapper  ark                       buildhash
\end{lstlisting}
\end{latin}

در این دستور، گزینه‌ی \cmd{-3} متن را در سه ستون سازمان‌دهی می‌کند و \cmd{-w 65} عرض صفحه را به ۶۵ نویسه محدود می‌نماید تا با ابعاد استاندارد نمایشگر یا کاغذ مطابقت داشته باشد. خروجی شامل یک سرآیند خودکار است که تاریخ، زمان و شماره‌ی صفحه را نمایش می‌دهد.

\section{ارسال عملیات چاپ به چاپگر}

سیستم چاپ \en{CUPS} از دو پروتکل استاندارد که به‌طور تاریخی در سیستم‌های شبه‌یونیکس برای اجرای عملیات چاپ استفاده می‌شده‌اند، پشتیبانی می‌کند. رویکرد نخست که به مدل \en{Berkeley} یا \en{LPD} شهرت دارد (و ریشه در توزیع نرم‌افزاری برکلی یا همان \en{BSD} دارد)، از برنامه‌ی \cmd{lpr} استفاده می‌کند. رویکرد دوم که با نام \en{SysV} شناخته می‌شود (برگرفته از نسخه \en{System V} یونیکس)، ابزار \cmd{lp} را به کار می‌گیرد. هر دو ابزار عملکردی تقریباً همسان دارند و انتخاب میان آن‌ها عمدتاً تابع عادت یا ترجیح شخصی کاربر است.

\subsection{دستور lpr: چاپ فایل‌ها (شیوه برکلی)}

برنامه‌ی \cmd{lpr} جهت ارسال فایل‌ها به چاپگر مورد استفاده قرار می‌گیرد. این ابزار به دلیل پذیرش ورودی استاندارد (\en{stdin})، به‌خوبی در پایپ‌لاین‌ها (\en{Pipelines}) قابل ادغام است. به عنوان نمونه، برای چاپ خروجیِ چندستونه‌ی فهرست محتویات دایرکتوری که پیش‌تر ایجاد کردیم، می‌توان از دستور زیر استفاده کرد:

\begin{latin}
\begin{lstlisting}
[me@linuxbox ~]$ ls /usr/bin | pr -3 | lpr
\end{lstlisting}
\end{latin}

در این حالت، گزارش مستقیماً به چاپگر پیش‌فرض سیستم ارسال می‌شود. جهت هدایت خروجی به چاپگری متفاوت، از گزینه‌ی \cmd{-P} به شیوه‌ی زیر استفاده می‌شود:

\begin{latin}
\begin{lstlisting}
lpr -P printer_name
\end{lstlisting}
\end{latin}

در دستور فوق، \file{printer\_name} نشان‌دهنده‌ی نام چاپگر مقصد است. برای مشاهده‌ی فهرست کامل چاپگرهای شناسایی‌شده در سیستم، می‌توان دستور زیر را به کار برد:

\begin{latin}
\begin{lstlisting}
[me@linuxbox ~]$ lpstat -a
\end{lstlisting}
\end{latin}

\begin{note}
\textbf{نکته:} بسیاری از توزیع‌های لینوکس قابلیت تعریف یک «چاپگر مجازی» را دارند که خروجی را به‌جای چاپ بر روی کاغذ، در قالب فایل \en{PDF} ذخیره می‌کند. این ویژگی ابزاری کارآمد برای آزمایش صحت دستورات چاپ بدون اتلاف منابع فیزیکی است. برای اطمینان از وجود این قابلیت، تنظیمات سیستم چاپ خود را بررسی کنید؛ در برخی توزیع‌ها، فعال‌سازی این امکان مستلزم نصب بسته‌های تکمیلی نظیر \cmd{cups-pdf} است.
\end{note}

جدول زیر پرکاربردترین گزینه‌های ابزار \cmd{lpr} را شرح می‌دهد:

\begin{table}[htbp]
\centering
\caption{گزینه‌های متداول دستور \cmd{lpr}}
\label{tab:lpr-options}
\begin{tabularx}{\textwidth}{l X}
\toprule
\textbf{گزینه} & \textbf{شرح عملکرد} \\
\midrule
\cmd{-\# number} & تعیین تعداد نسخه‌های چاپی بر اساس عدد مشخص شده. \\
\addlinespace
\cmd{-p} & فعال‌سازی قالب چاپ آراسته (\en{Pretty Print})؛ این گزینه در پیاده‌سازی \en{BSD} از دستور \cmd{pr} وجود دارد و باعث می‌شود هر صفحه با یک سرآیند سایه‌دار شامل تاریخ، زمان، عنوان سند و شماره صفحه چاپ شود. \\
\addlinespace
\cmd{-P printer} & تعیین نام چاپگر مورد نظر جهت خروجی. در صورت عدم استفاده از این گزینه، خروجی به چاپگر پیش‌فرض سیستم ارسال می‌گردد. \\
\addlinespace
\cmd{-r} & حذف خودکار فایل‌ها بلافاصله پس از اتمام فرآیند چاپ. این گزینه به‌ویژه برای مدیریت فایل‌های موقتی (\en{Temporary Files}) ایجاد شده توسط برنامه‌ها مفید است. \\
\bottomrule
\end{tabularx}
\end{table}

\begin{note}
\textbf{نکته:} این گزینه در نسخهٔ \en{GNU pr} (که در اکثر توزیع‌های لینوکس به کار می‌رود) پشتیبانی نمی‌شود.
\end{note}

\subsection{دستور lp: چاپ فایل‌ها (شیوه System V)}

برنامه‌ی \cmd{lp} نیز همانند \cmd{lpr} برای ارسال فایل‌ها یا ورودی استاندارد به چاپگر طراحی شده است. تفاوت اصلی این ابزار با نسخه \en{Berkeley} در پشتیبانی از مجموعه‌ای متفاوت و تا حدودی پیشرفته‌تر از گزینه‌های کنترلی است. جدول زیر کاربردی‌ترین گزینه‌های دستور \cmd{lp} را شرح می‌دهد:

\begin{table}[htbp]
\centering
\caption{گزینه‌های متداول دستور \cmd{lp}}
\label{tab:lp-options}
\begin{tabularx}{\textwidth}{l X}
\toprule
\textbf{گزینه} & \textbf{شرح عملکرد} \\
\midrule
\cmd{-d printer} & تعیین چاپگر مقصد. در صورت عدم استفاده از این گزینه، چاپگر پیش‌فرض سیستم انتخاب می‌شود. \\
\addlinespace
\cmd{-n number} & تعیین تعداد نسخه‌های چاپی. \\
\addlinespace
\cmd{-o landscape} & تنظیم جهت‌گیری صفحه به‌صورت افقی (\en{Landscape}). \\
\addlinespace
\cmd{-o fitplot} & تغییر مقیاس فایل جهت انطباق کامل با ابعاد کاغذ. این گزینه بیشتر برای چاپ تصاویر (مانند \en{JPEG}) کاربرد دارد. \\
\addlinespace
\cmd{-o scaling=number} & مقیاس‌بندی محتوا بر حسب درصد. عدد ۱۰۰ معادل اندازهٔ کامل صفحه است؛ مقادیر کمتر، کوچک‌نمایی و مقادیر بیشتر، چاپ را در چندین صفحه پخش می‌کنند. \\
\addlinespace
\cmd{-o cpi=number} & تنظیم تعداد کاراکتر در هر اینچ (\en{Characters Per Inch}) برای فایل‌های متنی. مقدار پیش‌فرض ۱۰ است. \\
\addlinespace
\cmd{-o lpi=number} & تنظیم تعداد خطوط در هر اینچ (\en{Lines Per Inch}) برای فایل‌های متنی. مقدار پیش‌فرض ۶ است. \\
\addlinespace
\cmd{-o page-bottom=points} & تنظیم حاشیهٔ پایین صفحه بر حسب پوینت (هر ۷۲ پوینت = ۱ اینچ). \\
\addlinespace
\cmd{-o page-left=points} & تنظیم حاشیهٔ چپ صفحه بر حسب پوینت. \\
\addlinespace
\cmd{-o page-right=points} & تنظیم حاشیهٔ راست صفحه بر حسب پوینت. \\
\addlinespace
\cmd{-o page-top=points} & تنظیم حاشیهٔ بالای صفحه بر حسب پوینت. \\
\addlinespace
\cmd{-P pages} & تعیین شماره صفحاتی که باید چاپ شوند. می‌توان از لیست‌های جدا شده با کاما یا بازه‌های عددی استفاده کرد (مانند \cmd{1-10,3,5,7}). \\
\bottomrule
\end{tabularx}
\end{table}

باید به این نکته توجه داشت که بسیاری از آرگومان‌های دستوری در \cmd{lp} (نظیر تنظیمات حاشیه، تراکم نویسه در اینچ و تعداد خطوط در اینچ) توسط فیلترهای داخلی سیستم چاپ مدیریت می‌شوند؛ لذا این گزینه‌ها لزوماً بر روی تمامی قالب‌های فایلی (به‌ویژه فایل‌های گرافیکی پیچیده) رفتار یکسانی ندارند. این تنظیمات بیشترین کارایی و دقت خود را هنگام پردازش فایل‌های متنی ساده (\en{Plain Text}) نشان می‌دهند.

در مثال زیر، فهرست محتویات دایرکتوری را مجدداً تولید می‌کنیم؛ با این تفاوت که تراکم خروجی را بر روی ۱۲ نویسه در اینچ (\en{CPI}) و ۸ سطر در اینچ (\en{LPI}) تنظیم کرده و حاشیه‌ی چپ را معادل نیم اینچ در نظر می‌گیریم. توجه داشته باشید که پارامترهای ابزار \cmd{pr} نیز باید متناسب با ابعاد جدید صفحه کالیبره شوند:

\begin{latin}
\begin{lstlisting}
[me@linuxbox ~]$ ls /usr/bin | pr -4 -w 90 -l 88 | lp -o page-left=36 -o cpi=12 -o lpi=8
\end{lstlisting}
\end{latin}

این پایپ‌لاین (\en{Pipeline})، فهرستی چهار ستونه با قلمی کوچک‌تر از حالت پیش‌فرض ایجاد می‌کند. افزایش شاخص نویسه در هر اینچ (\en{CPI}) این امکان را فراهم می‌آورد تا تعداد ستون‌های بیشتری در عرض صفحه گنجانده شود.

\subsection{دستور a2ps: تبدیل و آماده‌سازی خروجی}

برنامه‌ی \cmd{a2ps} (که در اکثر مخازن رسمی توزیع‌های لینوکس در دسترس است) ابزاری بسیار کارآمد محسوب می‌شود. همان‌طور که از نام آن استنباط می‌شود، این برنامه در اصل یک مبدل فرمت است، اما قابلیت‌هایی فراتر از یک تبدیل ساده ارائه می‌دهد. عنوان این ابزار در ابتدا مخفف «\en{ASCII to PostScript}» بود و جهت آماده‌سازی فایل‌های متنی برای چاپگرهای \en{PostScript} طراحی شده بود؛ اما با تکامل آن در طول سال‌ها، اکنون به معنای «\en{Anything to PostScript}» شناخته می‌شود.

اگرچه نام این برنامه به تبدیل فرمت اشاره دارد، اما در عمل یک راهکار جامع برای چاپ به شمار می‌رود. خروجی پیش‌فرض آن به‌جای هدایت به خروجی استاندارد (\en{stdout})، مستقیماً به چاپگر پیش‌فرض سیستم ارسال می‌گردد. رفتار استاندارد \cmd{a2ps} بر پایه‌ی «چاپ آراسته» (\en{Pretty Print}) بنا شده است، به این معنا که به‌طور خودکار جنبه‌های بصری خروجی را بهبود می‌بخشد.

در مثال زیر، از این برنامه برای تولید یک فایل \en{PostScript} بر روی دسکتاپ استفاده می‌کنیم:

\begin{latin}
\begin{lstlisting}
[me@linuxbox ~]$ ls /usr/bin | pr -3 -t | a2ps -o ~/Desktop/ls.ps -L 66
[stdin (plain): 11 pages on 6 sheets]
[Total: 11 pages on 6 sheets] saved into the file `/home/me/Desktop/ls.ps'
\end{lstlisting}
\end{latin}

در این مرحله، جریان داده‌ها ابتدا توسط ابزار \cmd{pr} و با به‌کارگیری گزینه‌ی \cmd{-t} (جهت حذف سرآیند‌ها و پانویس‌ها) فیلتر می‌شود؛ سپس خروجی حاصل به \cmd{a2ps} هدایت می‌گردد. در این دستور، با استفاده از گزینه‌ی \cmd{-o} مسیر ذخیره‌سازی فایل خروجی تعیین شده و با پارامتر \cmd{-L 66} (تنظیم ۶۶ سطر در هر صفحه)، تطابق کامل با ساختار صفحه‌بندیِ \cmd{pr} برقرار می‌شود. چنانچه فایل تولیدی با نرم‌افزار مناسبی باز شود، خروجی مشابه آنچه در شکل ۶ آمده است، نمایش داده خواهد شد.

\bookfigure{img-06.png}{خروجی نهایی دستور \cmd{a2ps}}

همان‌طور که ملاحظه می‌شود، چیدمان پیش‌فرض خروجی بر اساس الگوی \en{``two up''} (درج دو صفحه در یک برگ) تنظیم شده است که موجب می‌شود محتوای دو صفحه‌ی منطقی بر روی یک سطح فیزیکی چاپ گردد. علاوه بر این، \cmd{a2ps} به‌صورت خودکار سرآیندها و پانویس‌های شکیلی را جهت ارتقای بصری اسناد اعمال می‌کند.

ابزار \cmd{a2ps} از طیف گسترده‌ای از تنظیمات پشتیبانی می‌کند که گزینه‌های پرکاربرد آن در جدول زیر خلاصه شده است:

\begin{longtable}{>{\raggedright\arraybackslash}p{0.32\textwidth} p{0.60\textwidth}}
\caption{گزینه‌های کاربردی دستور \cmd{a2ps}}
\label{tab:a2ps-options} \\
\toprule
\textbf{گزینه} & \textbf{شرح عملکرد} \\
\midrule
\endfirsthead

\multicolumn{2}{c}{\tablename\ \thetable{} (ادامه) \rule[-0.4em]{0pt}{1.4em}} \\
\toprule
\textbf{گزینه} & \textbf{شرح عملکرد} \\
\midrule
\endhead

\midrule
\multicolumn{2}{r}{\footnotesize ادامه در صفحه‌ی بعد \ldots} \\
\endfoot

\bottomrule
\endlastfoot

\cmd{--center-title=text} & تنظیم متن بخش میانی سرآیند. \\
\addlinespace
\cmd{--columns=number} & تعیین تعداد ستون‌های چیدمان در هر برگ (مقدار پیش‌فرض: ۲). \\
\addlinespace
\cmd{--footer=text} & تنظیم متن پاورقی صفحه. \\
\addlinespace
\cmd{--guess} & تشخیص و گزارش نوع فایل‌های ورودی (مفید برای پیش‌بینی نحوه قالب‌بندی فایل توسط برنامه). \\
\addlinespace
\cmd{--left-footer=text} & تنظیم پاورقی سمت چپ. \\
\addlinespace
\cmd{--left-title=text} & تنظیم عنوان سمت چپ در سرآیند. \\
\addlinespace
\cmd{--line-numbers=interval} & شماره‌گذاری خطوط خروجی با فواصل عددی مشخص (\en{interval}). \\
\addlinespace
\cmd{--list=defaults} & نمایش پیکربندی و تنظیمات پیش‌فرض برنامه. \\
\addlinespace
\cmd{--pages=range} & محدود کردن فرآیند چاپ به بازه‌ی صفحات تعیین‌شده. \\
\addlinespace
\cmd{--right-footer=text} & تنظیم پاورقی سمت راست. \\
\addlinespace
\cmd{--right-title=text} & تنظیم عنوان سمت راست در سرآیند. \\
\addlinespace
\cmd{--rows=number} & تعیین تعداد ردیف‌های چیدمان صفحات در هر برگ (مقدار پیش‌فرض: ۱). \\
\addlinespace
\cmd{-B} & حذف کامل سرآیندها از خروجی. \\
\addlinespace
\cmd{-b text} & جایگزینی متن سرآیند با عبارت دلخواه. \\
\addlinespace
\cmd{-f size} & تنظیم اندازه‌ی قلم بر حسب واحد پوینت (\en{point}). \\
\addlinespace
\cmd{-l number} & تعیین تعداد کاراکتر در هر سطر (مناسب برای همگام‌سازی صفحه‌بندی با ابزارهایی نظیر \cmd{pr}). \\
\addlinespace
\cmd{-L number} & تعیین تعداد سطرها در هر صفحه. \\
\addlinespace
\cmd{-M name} & مشخص کردن نوع کاغذ مصرفی (مانند \en{A4} یا \en{Letter}). \\
\addlinespace
\cmd{-n number} & تعیین تعداد نسخه‌های چاپی از هر صفحه. \\
\addlinespace
\cmd{-o file} & هدایت خروجی به یک فایل مشخص (استفاده از \cmd{-} خروجی را به \en{stdout} می‌فرستد). \\
\addlinespace
\cmd{-P printer} & تعیین چاپگر مقصد؛ در صورت عدم انتخاب، از چاپگر پیش‌فرض استفاده می‌شود. \\
\addlinespace
\cmd{-R} & انتخاب حالت چاپ عمودی (\en{Portrait}). \\
\addlinespace
\cmd{-r} & انتخاب حالت چاپ افقی (\en{Landscape}). \\
\addlinespace
\cmd{-T number} & تعیین فاصله‌ی نقاط توقف \en{Tab} بر حسب تعداد نویسه؛ به عبارت دیگر، موقعیت‌های افقی در هر سطر را مشخص می‌کند که مکان‌نما پس از فشردن کلید \en{Tab} به آن‌ها منتقل می‌شود. \\
\addlinespace
\cmd{-u text} & درج متن مشخص‌شده به‌عنوان پس‌زمینه (\en{Watermark}) در تمام صفحات. \\

\end{longtable}

موارد فوق تنها بخشی از قابلیت‌های این ابزار قدرتمند است؛ \cmd{a2ps} برای مدیریت دقیق‌تر خروجی، گزینه‌های تخصصی بسیار بیشتری را در اختیار کاربران قرار می‌دهد.

\begin{note}
\textbf{نکته:} ابزار دیگری نیز با نام \cmd{enscript} برای تبدیل متون به قالب \en{PostScript} وجود دارد که بخش بزرگی از قابلیت‌های آرایشی و تکنیک‌های چاپ \cmd{a2ps} را پوشش می‌دهد؛ با این تفاوت که برخلاف \cmd{a2ps} که انواع مختلف داده را پردازش می‌کند، \cmd{enscript} صرفاً پذیرای ورودی‌های متنی است.
\end{note}

\section{نظارت و مدیریت عملیات چاپ}

معماری سیستم‌های چاپ در یونیکس به‌گونه‌ای پی‌ریزی شده است که توانایی مدیریت هم‌زمان چندین درخواست چاپ از سوی کاربران مختلف را داشته باشند؛ \en{CUPS} نیز دقیقاً بر پایه همین منطق طراحی شده است. برای هر چاپگر، یک صف چاپ (\en{Print Queue}) تعریف می‌شود که وظایف تا زمان آزادسازی منابع و ارسال به سخت‌افزار، در آن نگهداری می‌شوند. \en{CUPS} مجموعه‌ای از ابزارهای خط فرمان را برای نظارت بر وضعیت چاپگرها و مدیریت این صف‌ها در اختیار کاربر قرار می‌دهد. همانند برنامه‌های \cmd{lpr} و \cmd{lp}، این ابزارهای مدیریتی نیز بر اساس استانداردهای کلاسیک سیستم‌های \en{Berkeley} و \en{System V} مدل‌سازی شده‌اند.

\subsection{دستور lpstat: پایش وضعیت سیستم چاپ}

فرمان \cmd{lpstat} جهت شناسایی نام چاپگرها و بررسی وضعیت دردسترس‌بودن آن‌ها در سیستم به کار می‌رود. به‌طور مثال، در سیستمی که مجهز به یک چاپگر فیزیکی (با شناسه \en{``printer''}) و یک چاپگر مجازی \en{PDF} (با شناسه \en{``PDF''}) است، وضعیت پذیرش درخواست‌ها بدین صورت استخراج می‌شود:

\begin{latin}
\begin{lstlisting}
[me@linuxbox ~]$ lpstat -a
PDF accepting requests since Mon 08 Dec 2024 03:05:59 PM EST
printer accepting requests since Tue 24 Feb 2025 08:43:22 AM EST
\end{lstlisting}
\end{latin}

افزون بر این، برای دستیابی به جزئیات دقیق‌تر از پیکربندی عملیاتیِ زیرسیستم چاپ، می‌توان از گزینه‌ی \cmd{-s} استفاده کرد:

\begin{latin}
\begin{lstlisting}
[me@linuxbox ~]$ lpstat -s
system default destination: printer
device for PDF: cups-pdf:/
device for printer: ipp://print-server:631/printers/printer
\end{lstlisting}
\end{latin}

در خروجی فوق، مشخص است که \en{``printer''} به عنوان مقصد پیش‌فرض سیستم (\en{Default Destination}) تعیین شده است. همچنین مشاهده می‌کنیم که این چاپگر از طریق شبکه و با بهره‌گیری از پروتکل چاپ اینترنتی (\en{IPP}) به سرویس‌دهنده چاپ (\en{Print Server}) متصل شده است.

جدول زیر پرکاربردترین گزینه‌های ابزار \cmd{lpstat} را شرح می‌دهد:

\begin{table}[htbp]
\centering
\caption{گزینه‌های متداول دستور \cmd{lpstat}}
\label{tab:lpstat-options}
\begin{tabularx}{\textwidth}{l X}
\toprule
\textbf{گزینه} & \textbf{شرح عملکرد} \\
\midrule
\cmd{-a [printer...]} & نمایش وضعیت پذیرش درخواست در صف انتظار چاپگرهای مشخص‌شده. این گزینه صرفاً ناظر بر آمادگی صف برای دریافت کارهای جدید است و بیانگر وضعیت عملیاتیِ خودِ چاپگر فیزیکی نیست. در صورت عدم تعیین نام، وضعیت تمامی صف‌ها گزارش می‌شود. \\
\addlinespace
\cmd{-d} & نمایش نام چاپگر پیش‌فرض سیستم. \\
\addlinespace
\cmd{-p [printer...]} & نمایش وضعیت کنونی و آمادگی عملیاتیِ چاپگرهای مشخص‌شده. در صورت عدم تعیین نام، وضعیت تمامی چاپگرها گزارش می‌شود. \\
\addlinespace
\cmd{-r} & نمایش وضعیتِ سرویس‌دهنده (\en{Server})، یعنی همان پردازش پس‌زمینه‌ی (\en{Daemon}) سیستم چاپ. \\
\addlinespace
\cmd{-s} & ارائه‌ی گزارش خلاصه از وضعیت کلی سیستم چاپ، شامل مقاصد پیش‌فرض و آدرس دستگاه‌ها. \\
\addlinespace
\cmd{-t} & ارائه‌ی گزارشی جامع و کامل از تمامی جزئیاتِ وضعیت سیستم چاپ. \\
\bottomrule
\end{tabularx}
\end{table}

\subsection{دستور lpq: پایش صف انتظار چاپ}

جهت بازبینی وضعیت جاری صف چاپ (\en{Print Queue}) و رهگیری وظایف موجود در آن، از برنامه‌ی \cmd{lpq} استفاده می‌شود. این ابزار امکان پایش عملیات‌های در حال اجرا و کارهای معلق را فراهم می‌کند. در نمونه‌ی زیر، وضعیت یک صف خالی برای چاپگر پیش‌فرض سیستم (با نام \en{``printer''}) نمایش داده شده است:

\begin{latin}
\begin{lstlisting}
[me@linuxbox ~]$ lpq
printer is ready
no entries
\end{lstlisting}
\end{latin}

چنانچه نام چاپگر به‌طور صریح (با استفاده از گزینه‌ی \cmd{-P}) تعیین نشود، وضعیت چاپگر پیش‌فرض سیستم گزارش می‌گردد. در صورتی که درخواستی را به چاپگر ارسال و سپس بلافاصله صف را بررسی کنیم، آن مورد در فهرست کارهای جاری مشاهده خواهد شد:

\begin{latin}
\begin{lstlisting}
[me@linuxbox ~]$ ls *.txt | pr -3 | lp
request id is printer-603 (1 file(s))
[me@linuxbox ~]$ lpq
printer is ready and printing
Rank    Owner   Job     File(s)               Total Size
active  me      603     (stdin)               1024 bytes
\end{lstlisting}
\end{latin}

در خروجی فوق، شناسهٔ اختصاصی کار (\en{Job ID}) و وضعیت عملیاتی آن (مانند \en{active}) به‌منظور مدیریت و نظارت بر وظایف چاپ در مراحل بعدی، نمایش داده شده است.

\subsection{دستورات lprm و cancel: ابطال وظایف چاپ}

سامانه‌ی \en{CUPS} دو برنامه‌ی مجزا را جهت متوقف‌سازی فرآیندهای چاپ و حذف آن‌ها از صف انتظار ارائه می‌دهد؛ ابزار \cmd{lprm} که مبتنی بر استاندارد \en{Berkeley} است و ابزار \cmd{cancel} که از شیوه‌ی \en{System V} پیروی می‌کند. اگرچه این دو ابزار در پارامترهای عملیاتی تفاوت‌های جزئی دارند، اما کارکرد بنیادین آن‌ها یکسان است. با ارجاع به مثال قبلی، می‌توانیم وظیفه‌ی چاپ ارسالی را به‌صورت زیر متوقف و ابطال کنیم:

\begin{latin}
\begin{lstlisting}
[me@linuxbox ~]$ cancel 603
[me@linuxbox ~]$ lpq
printer is ready
no entries
\end{lstlisting}
\end{latin}

هر یک از این دستورات دارای گزینه‌های اختصاصی جهت حذف یک‌باره‌ی تمامی کارهای متعلق به یک کاربر خاص، پاکسازی صف یک چاپگر معین یا ابطال چندین شناسهٔ کاری به‌صورت هم‌زمان هستند. جزئیات فنی و پارامترهای تکمیلی این ابزارها در صفحات راهنما (\en{man pages}) به‌دقت مستند شده است.

\section{جمع‌بندی}

در این فصل، به بررسی جامع سیستم چاپ در محیط‌های شبه‌یونیکس پرداختیم. مشاهده کردیم که چگونه معماری چاپگرهای کلاسیک — از چاپگرهای ضربه‌ای با قلم‌های تک‌فاصله تا چاپگرهای گرافیکی مدرن با زبان‌های توصیف صفحه مانند \en{PostScript} — بر طراحی زیرسیستم‌های چاپ تأثیر گذاشته است. ابزارهای متنوعی مانند \cmd{pr} برای آماده‌سازی متن، \cmd{lpr} و \cmd{lp} برای ارسال وظایف به چاپگر، و \cmd{a2ps} برای تبدیل انواع داده به قالب \en{PostScript} معرفی شدند. همچنین با دستورات مدیریت صف مانند \cmd{lpstat}، \cmd{lpq}، \cmd{lprm} و \cmd{cancel} آشنا شدیم که امکان نظارت و کنترل دقیق بر وظایف چاپ را فراهم می‌کنند. در نهایت، مشاهده کردیم که رابط خط فرمان با ارائه‌ی گزینه‌های متعدد، کنترل بی‌نظیری بر زمان‌بندی، مدیریت صف‌های انتظار و تنظیمات خروجی در اختیار کاربر قرار می‌دهد.

\section{منابع مطالعه تکمیلی}

مقاله جامعی پیرامون زبان توصیف صفحه \en{PostScript}:

\begin{latin}
\url{http://en.wikipedia.org/wiki/PostScript}
\end{latin}

سامانه چاپ مشترک یونیکس (\en{Common Unix Printing System - CUPS}):

\begin{latin}
\url{http://en.wikipedia.org/wiki/Common_Unix_Printing_System}\\
\url{https://openprinting.github.io/cups/}
\end{latin}

سیستم‌های چاپ برکلی (\en{Berkeley}) و \en{System V}:

\begin{latin}
\url{http://en.wikipedia.org/wiki/Berkeley_printing_system}\\
\url{http://en.wikipedia.org/wiki/System_V_printing_system}
\end{latin}